\section{A finite Sperner lemma with its incidence proof}
\label{sec:finite-sperner}

Fix integers $N\geq 1$ and $d\geq 0$. Use cumulative coordinates for the
simplex grid,
\[
 G_d(N)=\{u\in\{0,\ldots,N\}^d: u_1\leq\cdots\leq u_d\}.
\]
Set $u_0=0$, $u_{d+1}=N$ and $a_i(u)=u_{i+1}-u_i$ for $0\leq i\leq d$.
The barycentric point represented by $u$ is $(a_0(u),\ldots,a_d(u))/N$.
The grid $G_0(N)$ consists of the empty tuple and its single gap is $N$.
A coloring is a table $c$ with values in $\{0,\ldots,d\}$ at grid points.
Its boundary condition is
\[
                       a_{c(u)}(u)>0\qquad(u\in G_d(N)).
\]

\paragraph{The cells.}
For $d>0$, take a sorted base $b\in\{0,\ldots,N-1\}^d$ and a permutation
$p_1,\ldots,p_d$ of $1,\ldots,d$. Its vertices are
\[
        v^0=b,\qquad v^s=b+\sum_{t=1}^s e_{p_t}
                     \quad(1\leq s\leq d).
\]
Keep this cell when all its vertices are sorted. Equivalently, within
each block of equal entries of $b$, the coordinates must be incremented
in decreasing index order. Indeed, incrementing the left of two equal
neighbors first violates their order. Two neighbors starting at distinct
integers stay ordered after either one is incremented once.
For $d=0$ use the single one-vertex cell. For any $s\leq t$ the
coordinates of $v^t-v^s$ belong to $\{0,1\}$. Extending this difference
by zero at the two endpoints shows
\begin{equation}\label{eq:sperner-gap-mesh}
              |a_i(v^t)-a_i(v^s)|\leq 1\qquad(0\leq i\leq d).
\end{equation}

The cells admit an enumeration with exactly $N^d$ entries. For a word
$w=(w_1,\ldots,w_d)\in\{0,\ldots,N-1\}^d$, sort it to obtain $b$.
Assign successive occurrences of each value in $w$ to the indices of its
block in $b$ from right to left. The index assigned to position $t$ is
$p_t$. The inverse reads $w_t=b_{p_t}$, because the validity criterion
uniquely fixes the order inside each block. This enumeration also works
for $N=1$ and for the empty word.

\paragraph{Finite Sperner lemma.}
For every coloring satisfying the boundary condition, the number of
cells with all $d+1$ colors is odd. The first such cell in word order can
therefore be found by scanning at most $N^d$ words.

\paragraph{Counting incidences.}
Suppose $d>0$. An incidence $(C,k)$ consists of a cell and the face
obtained by deleting its vertex $v^k$, $0\leq k\leq d$. Count it if the
remaining $d$ vertices carry precisely the colors $0,\ldots,d-1$.
A cell missing one of these colors contributes zero. If it contains
them all, its one remaining color occurrence either is $d$ or repeats
one of these colors. In the first case just the deletion of the
$d$-colored vertex counts. In the second case precisely the two
deletions of the repeated color count. Thus the parity of the incidence
count is the parity of the fully colored cell count.

Here is an explicit pairing on the incidences. The formulas preserve
the ordered list of retained vertices and hence preserve the counted
condition.
\begin{enumerate}
\item If $0<k<d$, interchange $p_k$ and $p_{k+1}$, using the same
deletion index $k$. Only the deleted vertex changes. If the interchange
is valid, applying it twice returns to the original cell. If it is
invalid, let $\ell=p_{k+1}$ and $y=v^{k-1}$. Increasing coordinate
$\ell$ first must violate $y_\ell\leq y_{\ell+1}$, so $\ell<d$ and
$y_\ell=y_{\ell+1}$. There is no possible overflow because $\ell$
has not yet been incremented. Since the vertex after both moves is
sorted, the other increment must be $p_k=\ell+1$. The two coordinates
are equal at every retained vertex, as neither has been incremented
before these moves and both have been incremented afterwards. All
retained vertices have $a_\ell=0$, so the missing color $\ell<d$
prevents this incidence from counting.
\item If $k=0$, let $j=p_1$. Use base $b+e_j$, order
$(p_2,\ldots,p_d,j)$, and delete the last vertex. Its first $d$ vertices
are the retained vertices of $C$. For $j<d$ the validity of the
original first move gives $b_j<b_{j+1}$, so the new last vertex is
sorted and within the coordinate bounds. For $j=d$, the additional
move is valid precisely when $b_d<N-1$. Thus the sole unpaired case
has $j=d$, $b_d=N-1$. Every retained vertex then has $u_d=N$.
\item If $k=d$, let $j=p_d$. Use base $b-e_j$, order
$(j,p_1,\ldots,p_{d-1})$, and delete the first vertex. This operation
inverts the preceding rotation whenever both are defined. If $j>1$,
the original penultimate vertex gives $b_{j-1}+1\leq b_j$, so the new
base is nonnegative and sorted. For $j=1$ the construction fails
precisely when $b_1=0$. In that case all retained vertices have
$u_1=0$, so none carries color zero and the incidence does not count.
\end{enumerate}

Each defined pairing exchanges distinct incidences. For the middle
deletion it changes a permutation, and at the ends it exchanges
deletion indices $0$ and $d$. Its counted pairs therefore contribute an
even number. The remaining counted incidences are exactly those of the
second case on $u_d=N$. Deleting the final coordinate of their retained
vertices gives cells in $G_{d-1}(N)$. Conversely, append $N-1$ to a
lower-dimensional base, prepend the step $d$ to its permutation, and
delete the initial vertex. These constructions are inverse. On the
top face the coloring $u'\mapsto c(u',N)$ satisfies the lower-dimensional
boundary condition because the last gap vanishes and the earlier gaps
are exactly the gaps of $u'$. Consequently the boundary contributes
exactly one incidence per fully colored lower-dimensional cell.
In dimension zero the unique vertex has color zero. This proves the
lemma by parity descent through the top faces.

\paragraph{The induction formula in iPRA.}
To justify that last descent uniformly in the dimension, fix the
ambient dimension $D$, $N$ and the entire finite color table $T$ as
parameters. For $r\leq D$ use its restriction
\[
 c_r(u_1,\ldots,u_r)=T(u_1,\ldots,u_r,
                           \underbrace{N,\ldots,N}_{D-r}).
\]
The positive-gap test on $T$ implies it for every restriction,
since all appended gaps beyond index $r$ vanish. Let $P(r)$ be the
number of fully colored $r$-cells modulo two computed by the word
enumeration on $c_r$. For a fixed parameter table the statement
$P(\min(r,D))=1$ is an equation of primitive recursive terms.
The dimension-zero computation starts its quantifier-free induction.
For $r<D$ the preceding incidence calculation gives
$P(r+1)=P(r)$, and for $r\geq D$ the capped index is constant.

All finite tests here may be implemented as Boolean folds. For example,
$T$ can use $(N+1)^D$ address slots in base $D+2$, with arbitrary
values at unused addresses and the support and range tests restricted
to sorted tuples. Word enumeration, bounded sorting, prefix sums,
permutation rotations and color counts are primitive recursive.
At dimension $r$ the incidence loop has exactly $(r+1)N^r$ entries.
For a formal counting justification, order the paired indices
numerically and group by the smaller index of each pair. Its fiber
consists of exactly its two members, since the pairing is an involution
without fixed points. Summing the indicator of the grouping relation
over the two finite indices in either order gives twice the number of
smaller members. The boundary bijection is checked by the same
double-sum argument with singleton fibers. Interchanging such finite
sums follows by quantifier-free induction on their bounds, with
bounded prefix sums as terms. All grid membership and incidence
checks have explicit finite enumerations. In particular, the induction
formula has no quantifier over possible color tables.

